\section{Perspectiva de la industria: el avance de los modelos y los límites de la innovación en los Agent}

\subsection{Que «el modelo se vuelva más fuerte» y que «el Agent siga teniendo valor» no son contradictorios}

Ji Yichao (季逸超) rechaza la idea de que «si el modelo es suficientemente fuerte, la capa de Agent deja de tener valor». La razón es que las tareas reales dependen del estado del entorno externo, y el entorno no puede interiorizarse por completo en los parámetros. Aunque el modelo sea más fuerte, la orquestación de tareas, el control de permisos y la verificación de la ejecución siguen requiriendo que la capa de Agent los asuma.

\begin{importantbox}{Un juicio más sólido}
El avance de los modelos reduce parte del costo de implementación del Agent, pero no elimina la responsabilidad sistémica del Agent. La competencia futura se parece más a «capacidad del modelo multiplicada por capacidad del sistema» que a una disyuntiva entre uno u otro.
\end{importantbox}

\subsection{Prioridad multimodal: primero comprender la entrada, después embellecer la salida}

En la entrevista se menciona que, en muchos escenarios, el verdadero punto doloroso es «entender una entrada compleja», no «generar una salida atractiva». Por ejemplo, capturas de pantalla, gráficas, documentos semiestructurados o imágenes devueltas por herramientas: si el modelo distorsiona la comprensión en la etapa de entrada, es muy difícil corregirlo en la cadena posterior.

\begin{warningbox}{Errores comunes en lo multimodal}
Concentrar la atención en el extremo de generación (salida de imagen o video) produce con facilidad un efecto de demostración a corto plazo, pero aporta poco a la tasa de finalización de tareas del Agent. Para un producto de flujo de trabajo, el análisis de la entrada y el seguimiento del estado suelen ser más determinantes.
\end{warningbox}

\subsection{Consejo para quienes emprenden: primero definir el problema, después emparejar la tecnología}

Un consejo aplicable de la entrevista es: primero dejar claro el problema, y después decidir si se construye un modelo, un sistema o solo una capa de producto. Solo cuando la definición del problema es precisa la vía técnica puede converger. De lo contrario, el equipo terminará dando vueltas sin cesar en un mar de oportunidades que «parece que todas se pueden hacer».

\begin{knowledgebox}{Plantilla de definición del problema (lista para usarse)}
\begin{itemize}
    \item Objetivo del usuario: cuál es el entregable final
    \item Restricciones: tiempo, permisos, auditabilidad, límites de seguridad
    \item Criterios de éxito: tasa de finalización en un intento, costo de los reintentos, proporción de intervención humana
\end{itemize}
\end{knowledgebox}

\subsection{Resumen de la sección}

Para emprender en el área de los Agent, el optimismo debe apoyarse en un juicio estructurado: reconocer que el modelo seguirá avanzando y, al mismo tiempo, volver insustituible la capacidad de la capa de sistema.

